\chapter{Các công trình nghiên cứu liên quan}
\section{Phân đoạn ảnh và mô hình thị giác--ngôn ngữ}
[Giải thích semantic, instance và binary segmentation; giới thiệu vai trò của embedding ảnh--văn bản.]

\section{Phân loại phương pháp theo prompt}
\begin{table}[H]
\centering
\caption{Khung khảo sát các phương pháp liên quan. Điền nguồn và mô tả sau khi kiểm chứng.}
\label{tab:related}
\begin{tabularx}{\textwidth}{@{}l L L L@{}}
\toprule
Phương pháp & Loại prompt & Đầu ra & Cách dùng VLM / nguồn \\
\midrule
CLIPSeg & Text, ảnh mẫu & Mask nhị phân & [Bổ sung] \\
{[Phương pháp 2]} & [Bổ sung] & [Bổ sung] & [Bổ sung] \\
{[Phương pháp 3]} & [Bổ sung] & [Bổ sung] & [Bổ sung] \\
{[Phương pháp 4]} & [Bổ sung] & [Bổ sung] & [Bổ sung] \\
\bottomrule
\end{tabularx}
\end{table}

\section{Vị trí của CLIPSeg}
[So sánh loại đầu vào, độ linh hoạt và giới hạn đầu ra. Phân biệt số công bố từ paper với số tự đo trong Chương 4.]
